\section{Methodology}
\label{sec:methodology}

\subsection{Study Overview}
\label{sec:study_overview}

This study investigates the ability of Small Language Models (SLMs) to both generate correct code and identify errors in code generated by other models. The experiment combines code generation, automated classification, knowledge distillation, and specialized training to evaluate different dimensions of the problem.

Several language models generate solutions for standardized programming tasks. Each solution is executed against a test suite that objectively determines whether the code is correct or contains syntax, runtime, or functional errors. This ground truth classification is used both to evaluate the code generation capability of the models and as the basis for detection experiments.

The study includes a knowledge distillation process in which language models act as judges to generate natural language explanations for each identified error type. These explanations are associated with the reference labels and compose the data used for supervised training of a specialized error detection model via \textit{fine-tuning} with the QLoRA technique.

Detection capability is evaluated for both general-purpose models, without specific training, and the specialized model. In both cases, the model receives only the programming task description and the generated code, without access to reference labels or test execution results.

The detection task is treated as multilabel classification, since the same code may contain more than one type of error simultaneously.

\subsection{Evaluated Models}
\label{sec:evaluated_models}

The study evaluates language models with parameter counts ranging from 4B to 552B, as shown in Table~\ref{tab:models}. The table separates the small language models that are the main focus of the study from the two large models used only for the comparison with larger models.

\begin{table}[htbp]
\centering
\caption{Models evaluated in the study}
\label{tab:models}
\begin{tabular}{lc}
\toprule
\textbf{Model} & \textbf{Parameters} \\
\midrule
\multicolumn{2}{l}{\textit{Small language models}} \\
\midrule
Devstral Small 2 & \shortstack[c]{24B\\{\scriptsize\citep{rastogi2025devstral}}\\{\scriptsize\citep{mistralai2025devstral2}}} \\
GLM 4.7 Flash & \shortstack[c]{30B\\{\scriptsize\citep{glm45team2025glm45}}\\{\scriptsize\citep{zai2025glm47flash}}} \\
Gemma 3 4B & \shortstack[c]{4B\\{\scriptsize\citep{gemmateam2025gemma3}}} \\
Gemma 4 & \shortstack[c]{\{5.1, 8, 31\}B\\{\scriptsize\citep{gemmateam2026gemma4}}} \\
GPT-OSS 20B & \shortstack[c]{20B\\{\scriptsize\citep{openai2025gptoss}}} \\
Muse Glimmer 30B & \shortstack[c]{30B\\{\scriptsize\citep{meta2026museglimmer}}} \\
Ornith 1.5 35B & \shortstack[c]{35B\\{\scriptsize\citep{ornithteam2026ornith15}}} \\
Qwen 2.5 Coder & \shortstack[c]{\{7, 32\}B\\{\scriptsize\citep{hui2024qwen25coder}}} \\
Qwen 3 Coder 30B & \shortstack[c]{30B\\{\scriptsize\citep{yang2025qwen3}}\\{\scriptsize\citep{qwenteam2025qwen3coder}}} \\
Qwen 3.8 27B & \shortstack[c]{27B\\{\scriptsize\citep{qwenteam2026qwen38}}} \\
\midrule
\multicolumn{2}{l}{\textit{Large models (comparison)}} \\
\midrule
DeepSeek V4.1 Flash & \shortstack[c]{552B\\{\scriptsize\citep{deepseekai2026deepseekv41flash}}} \\
MiMo V2.5 & \shortstack[c]{310B\\{\scriptsize\citep{xiaomimimo2026mimov25}}} \\
\bottomrule
\end{tabular}
\end{table}

The local SLMs used in the initial code generation were executed with Q4\_K\_M quantization, which allowed larger models to be evaluated within the available computational budget. The Gemma models were also evaluated without precision reduction in the detection experiments. The fine-tuned models were executed without inference quantization; their adaptation used QLoRA with 4-bit NF4 quantization~\citep{dettmers2023qlora}, which is a training mechanism distinct from Q4\_K\_M inference quantization. The two large models were accessed through a hosted API, so no quantization level is assumed for them.

For the knowledge distillation process, three models were selected as judges: Muse Glimmer 30B, Qwen 3.8 27B, and Ornith 1.5 35B. These models were chosen because they are the largest general-purpose models and demonstrated strong performance in the code generation step.

For the specialized model training, the Gemma 4 E2B and E4B models were used as base models for \textit{fine-tuning}, along with Gemma 3 4B.

During the knowledge distillation and error detection evaluation, all models were configured with a temperature of 0.0 to ensure deterministic outputs. Across all experiments, both the SLMs and the larger comparison models were executed in direct-response mode, with thinking or reasoning mode and reasoning effort disabled whenever these controls were exposed by the model backend. Consequently, the models were not prompted to produce a chain-of-thought or a separate reasoning trace, and no such trace was retained or supplied as input to any subsequent phase.

\subsection{Programming Tasks and Code Generation}
\label{sec:programming_tasks}

The experiments utilize the EvalPlus benchmark~\citep{liu2023evalplus}, which extends both HumanEval~\citep{chen2021humaneval} and MBPP~\citep{austin2021mbpp} with additional test cases to increase the reliability of correctness assessment. EvalPlus provides 164 tasks from HumanEval and 399 tasks from MBPP, totaling 563 programming tasks.

Each task consists of a natural language problem description that specifies the required functionality. The problems cover diverse algorithmic challenges, including string manipulation, mathematical operations, and data structure operations. For HumanEval tasks, the model receives the function signature and must complete the function body. For MBPP tasks, the model generates the complete function implementation.

The code generation process uses a system prompt that instructs the model to act as a Python assistant that solves coding problems, returning only a single function implementation without tests, asserts, or example usage. The user prompt varies according to the benchmark: for HumanEval, it requests the completion of a function body; for MBPP, it requests the complete function implementation.

For the initial code generation, each model produced one solution per programming task, using a temperature of 1.0 to allow for diversity. The top-p and top-k parameters were kept at the defaults of each model or backend, and no fixed random seed was imposed.

The raw response is normalized before execution. The extraction step removes Markdown code fences when they wrap the whole response and detects whether the model already returned the entry-point function definition. For HumanEval completion tasks, the harness reuses the task prompt, which already contains the function header, and appends the normalized function body; when the model returned the complete function definition, that definition is used directly. For MBPP, the returned function definition is used when present; otherwise the required signature is prepended and the body is indented. This normalization does not remove arbitrary natural-language text that the model may produce before or after the code, so a response that mixes prose with code can still fail to compile. When extraction produces an incomplete source, for example when the function header is lost, the ground-truth compilation classifies the solution as Syntax.

\subsection{Ground Truth and Error Classification}
\label{sec:ground_truth}

The ground truth classification is determined by executing each generated solution against a custom test harness that evaluates the code at different levels of correctness. Labels are accumulated across the test cases, so a single solution may receive more than one label, with the exception of Correct.

The procedure starts with a compilation step. If the source cannot be compiled, it receives the Syntax label and the evaluation stops, so Syntax is exclusive and no other label is added. It is important to note that syntax errors may arise from different sources: (1) genuine Python syntax errors in the generated code, such as invalid import statements or malformed expressions; (2) formatting issues, such as the model returning natural language explanations or Markdown fences along with the code; or (3) extraction issues, such as the loss of the function header during code processing.

If the source compiles, it is executed in an isolated subprocess against the base test cases and the additional test cases. For each test case, if the execution raises an exception, the Runtime label is added; if the execution completes but returns a result different from the expected one, the Functional label is added. Runtime and Functional can therefore co-occur in the same solution, because some test cases may raise exceptions while others return incorrect results. A timeout, or a failure that prevents the execution of the candidate as a whole, is also labeled as Runtime.

The Correct label is assigned only when the source compiles, no test case raises an exception, and all test cases return the expected result. It is therefore exclusive: a solution labeled Correct cannot receive any other label, and a solution that receives any error label cannot receive Correct. The procedure is summarized in Figure~\ref{fig:ground_truth}.

\begin{figure}[htbp]
\centering
\begin{minipage}{0.95\columnwidth}
\begin{verbatim}
Input: generated code, EvalPlus tests
labels = {}
if code does not compile:
    labels += Syntax
    return labels
for each test case:
    execute code
    if exception or timeout:
        labels += Runtime
        continue
    if output != expected result:
        labels += Functional
if labels is empty:
    labels += Correct
return labels
\end{verbatim}
\end{minipage}
\caption{Ground truth labeling procedure. Runtime and Functional accumulate across test cases, whereas Syntax terminates the evaluation and Correct is assigned only when no other label is present.}
\label{fig:ground_truth}
\end{figure}

\subsection{Error Detection Experiment}
\label{sec:error_detection}

The error detection experiment evaluates whether language models can correctly identify errors in generated code without access to reference labels or test execution results. For each sample, the model receives only the programming task description and the generated code.

The models are prompted with a system prompt that defines the classification task. Table~\ref{tab:error_types} describes the four error categories used in this study.

\begin{table}[htbp]
\centering
\caption{Error categories and their definitions}
\label{tab:error_types}
% Keep the definitions table within a single column and let the text wrap.
\begin{tabular}{@{}p{0.24\columnwidth}@{}p{0.76\columnwidth}@{}}
\toprule
\textbf{Error Type} & \textbf{Definition} \\
\midrule
Syntax & When the provided code has a syntax error in the Python language, such as incorrect or incomplete code. \\
\midrule
Runtime & When the code is ``compilable'' but has an error when executed, such as a call to an invalid function, an undefined variable, etc. \\
\midrule
Functional & When the code is executable, but it has some deviation from the correct functioning of what was requested. \\
\midrule
Correct & When the code is well-formed and has no functional errors. \\
\bottomrule
\end{tabular}
\end{table}

The prompt instructs the model to return a JSON object containing a list of detected error types, each with a natural language explanation. The prompt emphasizes that the model should provide general explanations without mentioning specific test cases or input values.

All models evaluated in the code generation step are also evaluated in the error detection experiment. The models are evaluated under the same conditions, with a temperature of 0.0 to ensure deterministic outputs. The two large models used in the comparison with larger models follow exactly the same detection protocol as the SLMs: same set of examples, same task description, same input code, same system prompt, same class definitions, same output format, same metrics, and the same temperature. No additional information, tests, or ground truth is provided to them.

The Exact-set Accuracy measures the proportion of samples where the set of predicted error types exactly matches the set of ground truth error types. Additionally, Macro F1 and Macro Recall are computed by isolating each error class independently, providing a per-class analysis of detection capability.

\subsection{Specialized Model Training and Knowledge Distillation}
\label{sec:specialized_training}

This section describes both the knowledge distillation process used to generate training data and the specialized model training procedure.

\subsubsection{Knowledge Distillation}

The knowledge distillation process generates natural language explanations for each error type identified in the ground truth classification. For each generated code sample, a judge model receives the problem description, the generated code, the name of the model that generated the code, and the ground truth error labels. The judge is instructed to write a general explanation for each error type present, describing why the code can fail in that way without mentioning specific test cases or input values. The judge does not assign or correct the label; it explains a label that was already determined by execution.

Every example is processed by all three judges: Muse Glimmer 30B, Qwen 3.8 27B, and Ornith 1.5 35B. These models were chosen because they are the largest general-purpose models and demonstrated strong performance in the code generation step. Each judge produces its own rationale for the same code and the same reference label, and no voting or winner selection is performed. Each generated sample therefore contributes one training record per judge, that is, three records that share the code and the label but carry a different rationale. Each rationale associates an error type with the explanation produced by the judge. If a judge fails to produce a response for a sample, no record is written for that judge and the sample is retried in a later execution. The reference label always comes from the execution of the code, and the rationale is only the target text that accompanies it.

Training the specialized model on the label together with an explanation, rather than on the label alone, provides richer auxiliary supervision. Label-only supervision indicates which output is expected but offers little explicit information about the relation between the task requirement, the behavior of the code, and the error category, which in an imbalanced dataset can favor superficial correlations and shortcuts associated with the dominant class. Conditioning the target on an explanation encourages the model to associate its prediction with evidence in the code. This follows the rationale-based distillation of \citet{hsieh-etal-2023-distilling}, which uses rationales as additional supervision, and the explanation-based fine-tuning of \citet{ludan-etal-2023-explanation}, which reports reduced reliance on spurious cues. The explanations provide richer auxiliary supervision and are not claimed to eliminate bias or to guarantee robustness. In the quantitative evaluation, only the predicted classes enter the metrics; the textual quality of the explanations is not scored.

\subsubsection{Specialized Model Training}

The training dataset is the output of the knowledge distillation process described above. Each record in this dataset contains the problem description, the generated code, and the corresponding error labels with explanations generated by the judge models.

The base models selected for \textit{fine-tuning} are Gemma 4 E2B, Gemma 4 E4B, and Gemma 3 4B. Each training sample is formatted as a three-turn conversation. The system prompt describes the classification task and defines the four error categories. The user message contains the programming task description and the generated code. The assistant message contains the expected output: a JSON object with the detected error types and their explanations, formatted as a Markdown code block.

The training uses QLoRA (Quantized Low-Rank Adaptation) with 4-bit NF4 quantization~\citep{dettmers2023qlora}, applying low-rank adapters as introduced by LoRA~\citep{hu2022lora}. Table~\ref{tab:training_params} shows the hyperparameters explored during the search.

\begin{table}[htbp]
\centering
\caption{Training hyperparameters}
\label{tab:training_params}
\begin{tabular}{ll}
\toprule
\textbf{Parameter} & \textbf{Values} \\
\midrule
LoRA rank ($r$) & \{8, 16, 32\} \\
LoRA alpha ($\alpha$) & \{8, 16, 32\} \\
LoRA dropout & \{0.05, 0.1\} \\
Learning rate & \{5e-5, 1e-4\} \\
Bias & \{none, all, lora\_only\} \\
Epochs & \{3, 4\} \\
\bottomrule
\end{tabular}
\end{table}

The dataset is split into training, validation, and test sets. The split is performed at the task level, so all code samples for a given programming task are assigned to the same split, preventing data leakage where different solutions for the same task appear across splits. The validation set comprises 15\% of the data, and the test set comprises 20\%.

Four training configurations were evaluated. The first two use the original class distribution, without removing Syntax and without explicit balancing. The first is Accuracy-oriented: it uses the standard causal language-modeling loss and selects the best model by validation Accuracy. The second is Macro-F1-oriented: it uses the label-weighted loss, which assigns a higher weight to the class-name tokens, and selects the best model by validation Macro-F1.

The other two configurations balance the training split only. Correct examples are reduced, Syntax examples are removed from training, and the remaining examples are downsampled to a runtime:functional:correct ratio of 1:1:1 or 1:3:3, without duplicating rows. Validation and test remain unchanged, so balancing cannot move a task between splits. Both balanced configurations use the label-weighted loss; the third selects the best model by validation Accuracy, and the fourth by validation Macro-F1.

Hyperparameter search was performed with Optuna~\citep{akiba2019optuna} using the Tree-structured Parzen Estimator (TPE) sampler with multivariate sampling and constant-liar. The search space covers the base model, LoRA rank and alpha, LoRA dropout, learning rate, number of epochs, and bias configuration, as listed in Table~\ref{tab:training_params}. For each search, the best configuration is the one that maximizes the respective selection metric on the validation split, and the corresponding merged adapter is retained.

\subsection{Evaluation Metrics}
\label{sec:evaluation_metrics}

The evaluation uses metrics appropriate for multilabel classification scenarios. For each of the four error categories (syntax, runtime, functional, and correct), true positives (TP), false positives (FP), and false negatives (FN) are computed independently.

The per-class recall is calculated as $TP / (TP + FN)$, measuring the proportion of actual positive cases correctly identified for each error type. The per-class precision is calculated as $TP / (TP + FP)$, measuring the proportion of predicted positive cases that are actually correct.

The Macro F1 score is computed as the average of the F1 scores across all four error categories, where each category's F1 score is $2 \times TP / (2 \times TP + FP + FN)$. This metric treats all error types equally, regardless of their frequency in the dataset.

The Macro Recall is computed as the average of the recall scores across all four error categories, providing a measure of the overall detection capability across error types.

The Exact-set Accuracy measures the proportion of samples where the set of predicted error types exactly matches the set of ground truth error types. This strict metric requires a complete match of the multilabel prediction, reflecting the practical requirement that a detection system must correctly identify all present error types.

\subsection{Static Analysis Tools}
\label{sec:static_tools}

Two static analysis tools are used in this study, each exclusive to one failure category. Python compilation (\texttt{py\_compile}) is responsible for the syntax category, and Pyright is responsible for statically detectable runtime issues.

Python compilation, implemented with the language's built-in \texttt{compile} function~\citep{pythoncompile}, parses the source code into a code object and reports a syntax error when the source cannot be parsed. A compilation failure is mapped to Syntax, whereas a successful compilation is interpreted as the absence of a syntax signal. It is never used to detect Runtime or Functional errors, and it does not execute the code. Because it only parses the source, it is fast and deterministic.

Pyright is a standards-based static type checker for Python~\citep{microsoftpyright} that analyzes the source code without executing it. It builds a model of the program and reports conditions that would fail at execution. In this study it is invoked under the default type checking mode with \texttt{pyright --outputjson --level error --pythonversion 3.13}, so only diagnostics with error severity are parsed. A diagnostic is mapped to Runtime only when its rule belongs to the runtime-related set considered here, which includes \texttt{reportUndefinedVariable}, \texttt{reportUnboundVariable}, \texttt{reportPossiblyUnboundVariable}, \texttt{reportCallIssue}, \texttt{reportArgumentType}, \texttt{reportAssignmentType}, \texttt{reportAttributeAccessIssue}, \texttt{reportIndexIssue}, \texttt{reportOperatorIssue}, \texttt{reportReturnType}, \texttt{reportOptionalCall}, \texttt{reportOptionalMemberAccess}, \texttt{reportOptionalOperand}, \texttt{reportOptionalSubscript}, \texttt{reportOptionalIterable}, \texttt{reportOptionalContextManager}, \texttt{reportOptionalEnter}, \texttt{reportGeneralTypeIssues}, and \texttt{reportInvalidTypeArguments}. Diagnostics without a rule whose message indicates a parsing problem, such as an expected token or an indentation error, are mapped to Syntax. All other diagnostics, including style and type-annotation warnings, are ignored. Pyright detects runtime issues that compilation cannot, but it cannot detect functional errors, since these depend on the expected behavior of the program. If a tool fails to run or produces output that cannot be parsed, the failure is recorded as a tool error and no prediction is derived from it.

\subsection{Interactive Cycles with Tools}
\label{sec:interactive_cycles}

Two experiments use the static tools in their loop: one for code generation and one for code verification. Figures~\ref{fig:cycle_generation} and~\ref{fig:cycle_verification} summarize both flows.

\subsubsection{For Code Generation}

\begin{figure}[htbp]
\centering
\resizebox{0.6\columnwidth}{!}{%
\begin{tikzpicture}[
  font=\small,
  box/.style={draw, rounded corners=2pt, align=center, minimum width=3.1cm, minimum height=0.8cm, fill=blue!5},
  subtool/.style={draw, rounded corners=2pt, align=center, minimum width=1.5cm, minimum height=0.7cm, fill=orange!15},
  container/.style={draw, rounded corners=3pt, fill=orange!8, minimum width=4.6cm, minimum height=1.8cm},
  dec/.style={draw, diamond, aspect=2.2, align=center, inner sep=1pt, fill=green!10},
  term/.style={draw, rounded corners=2pt, align=center, minimum width=3.1cm, minimum height=0.8cm, fill=gray!12},
  arr/.style={-{Latex[length=2mm]}, thick},
  lbl/.style={font=\scriptsize}
]
\node[term] (g0) at (0,0) {Programming task};
\node[box]  (g1) [below=6mm of g0] {Model generates / revises code};
\node[container] (tf) [below=6mm of g1] {};
\node[font=\footnotesize] at ([yshift=-3mm]tf.north) {Tool feedback};
\node[subtool] at ([xshift=-1.05cm,yshift=-4.5mm]tf.center) {py\_compile};
\node[subtool] at ([xshift=1.05cm,yshift=-4.5mm]tf.center) {Pyright};
\node[box]  (g3) [below=6mm of tf] {Model self-review};
\node[dec]  (g4) [below=6mm of g3] {Submit?};
\node[term] (g5) [below=6mm of g4] {EvalPlus (once, final)};
\draw[arr] (g0)--(g1);
\draw[arr] (g1)--(tf);
\draw[arr] (tf)--(g3);
\draw[arr] (g3)--(g4);
\draw[arr] (g4)--node[lbl,right]{yes}(g5);
\draw[arr] (g4.east) -- ++(1.7,0) node[lbl,above,pos=0.5]{no: revise} |- (g1.east);
\end{tikzpicture}%
}
\caption{Code generation cycle. The model iterates between generation, static analysis, and self-review, and EvalPlus is applied only to the final candidate.}
\label{fig:cycle_generation}
\end{figure}

In the generation cycle, the model first produces an initial solution for a programming task. In each round, the candidate is analyzed by Python compilation and, only when it compiles, by Pyright, so the tools run in a short-circuit order and a candidate that does not compile does not reach Pyright. The resulting diagnostics are returned to the model, which performs a self-review and may either submit the current solution or return a revised implementation, following the broader paradigm of iterative self-feedback and self-debugging~\citep{madaan2023selfrefine,chen2024selfdebug}. The prompt contains only the task, the current candidate, and the current feedback; the full history of previous rounds is not provided. When no diagnostic is found, the feedback states that no syntax or statically detectable runtime issue was reported and still asks the model to assess the code, so the model may revise even without diagnostics. When the model revises, the returned implementation completely replaces the candidate, without any correctness validation before the substitution, and the cycle repeats, and when the model does not return the requested JSON, a code-only fallback is used. The cycle stops when the model submits or when the maximum of five rounds is reached, and both the generation and the review use a temperature of 0. The evaluation with EvalPlus is applied only once, to the final candidate, and its result is never included in any generation or self-review prompt, so the model has no access to the test outcome during the cycle. Because a revision is not guaranteed to be better than the previous candidate, the process is also analyzed in terms of improved, unchanged, and regressed outcomes, and regressions are retained rather than discarded.

\subsubsection{For Code Verification}

In the verification experiment, the model receives the programming task description and a held-out solution generated during the initial code generation, together with the static feedback produced by Python compilation and Pyright for that solution. This experiment is a controlled comparison: the same model, on the same sample, is evaluated with and without the static feedback. The version without feedback corresponds to the general error detection experiment, and the two versions share the same task, the same code, the same base prompt, the same temperature, and the same metrics, so only the static feedback differs. The feedback is presented as supporting evidence and never as a ground truth label. It is a short text that reports the compile status and, when the source compiles, the Pyright runtime status, followed by the diagnostics; when no issue is found, it states that no syntax or statically detectable runtime issue was reported and adds that this does not prove functional correctness. The prompt states that the tools can be incomplete, that the absence of a diagnostic does not imply that the code is correct, and that the tools do not detect functional errors in general. The model then returns the multilabel classification of the code into syntax, runtime, functional, and correct. The static feedback is used only as evidence for the classification and never replaces the reference labels, which are used exclusively afterward to compute the metrics.

\begin{figure}[htbp]
\centering
\resizebox{0.6\columnwidth}{!}{%
\begin{tikzpicture}[
  font=\small,
  box/.style={draw, rounded corners=2pt, align=center, minimum width=3.1cm, minimum height=0.8cm, fill=blue!5},
  subtool/.style={draw, rounded corners=2pt, align=center, minimum width=1.5cm, minimum height=0.7cm, fill=orange!15},
  container/.style={draw, rounded corners=3pt, fill=orange!8, minimum width=4.6cm, minimum height=1.8cm},
  term/.style={draw, rounded corners=2pt, align=center, minimum width=3.1cm, minimum height=0.8cm, fill=gray!12},
  arr/.style={-{Latex[length=2mm]}, thick}
]
\node[term] (v0) at (0,0) {Task + held-out code};
\node[container] (tf) [below=6mm of v0] {};
\node[font=\footnotesize] at ([yshift=-3mm]tf.north) {Tool feedback};
\node[subtool] at ([xshift=-1.05cm,yshift=-4.5mm]tf.center) {py\_compile};
\node[subtool] at ([xshift=1.05cm,yshift=-4.5mm]tf.center) {Pyright};
\node[box] (v3) [below=6mm of tf] {Model classifies};
\node[box] (v4) [below=6mm of v3] {Predicted labels};
\node[term] (v5) [below=6mm of v4] {Compare with ground truth};
\draw[arr] (v0)--(tf);
\draw[arr] (tf)--(v3);
\draw[arr] (v3)--(v4);
\draw[arr] (v4)--(v5);
\end{tikzpicture}%
}
\caption{Code verification cycle. The model classifies held-out code using the static feedback as supporting evidence.}
\label{fig:cycle_verification}
\end{figure}
